\FloatBarrier
\section{Qubit architectures and encodings}
\label{sec:qubit-architectures}

After selecting and retaining the $N=1$ charge state, we can keep
the electron in the lowest orbital--valley state $\ket{\chi_0}$ and
use its two spin states to store a qubit. With several confined
electrons, we can instead select two joint spin states and use them
as one logical qubit. The choice of this two-dimensional
\emph{computational subspace} is the encoding; transitions out of
it are leakage. Different encodings use the available physical
interactions differently: a field acts on an individual spin,
while an interaction between spins can act within a jointly
encoded space.

We retain one orbital--valley state per dot and assume predominantly
one electron per occupied dot.

\subsection{Single-spin qubits}
\label{sec:single-spin-qubits}

For one electron, we choose the logical basis
\begin{equation}
  \ket{0}=\ket{\chi_0}\otimes\ket{\downarrow},
  \qquad
  \ket{1}=\ket{\chi_0}\otimes\ket{\uparrow}.
  \label{eq:qubit-basis}
\end{equation}
The computational space is
$\Hilb_{\mathrm q}=\operatorname{span}\{\ket0,\ket1\}$,
and a pure qubit state is $\ket{\psi}=a\ket0+b\ket1$, with
$|a|^2+|b|^2=1$.
Leakage involves charge, orbital, or valley excitations.

Charge preparation alone does not fix the spin state.
For example, the two-dot $(1,1)$ charge configuration introduced
in Section~\ref{sec:multiple-quantum-dots}, with one orbital--valley
state per dot, contains four spin states:
\begin{equation}
  \ket{\downarrow\downarrow},\quad
  \ket{\downarrow\uparrow},\quad
  \ket{\uparrow\downarrow},\quad
  \ket{\uparrow\uparrow}.
  \label{eq:two-spin-basis}
\end{equation}
The arrows refer to dots 1 and 2. A charge stability diagram identifies
charge occupations and transitions, not the spin state within a charge
sector; its boundaries are derived in
Appendix~\ref{sec:measurement-double-dot-csd}.
The four states can represent two single-spin qubits or supply
a subspace for one encoded qubit~\cite{Burkard2023}.

\subsection{Singlet--triplet qubits}
\label{sec:singlet-triplet-qubits}

Two spins decompose as $\tfrac12\otimes\tfrac12=0\oplus1$.
We encode one qubit in the $m=0$ states: the singlet ($S=0$) and
the triplet ($S=1$),\footnote{The label $S$ in $\ket S$ stands for
``singlet''; the total-spin quantum number of this state is $S=0$.}
\begin{equation}
  \begin{aligned}
    \ket{0_{\mathrm{ST}}}=\ket{S}
      &=\frac{\ket{\uparrow\downarrow}-\ket{\downarrow\uparrow}}{\sqrt{2}},\\
    \ket{1_{\mathrm{ST}}}=\ket{T_0}
      &=\frac{\ket{\uparrow\downarrow}+\ket{\downarrow\uparrow}}{\sqrt{2}}.
  \end{aligned}
  \label{eq:st-basis}
\end{equation}
The remaining triplets
$\ket{T_+}=\ket{\uparrow\uparrow}$ and
$\ket{T_-}=\ket{\downarrow\downarrow}$ are spin-leakage states.

\subsection{Exchange-only qubits}
\label{sec:exchange-only-qubits}

Three spin-$1/2$ electrons can encode one logical qubit. Their spin
space decomposes as
\[
 \tfrac12\otimes\tfrac12\otimes\tfrac12
   =\tfrac12\oplus\tfrac12\oplus\tfrac32.
\]
Taking the $m=+1/2$ state from each of the two $S=1/2$
representations gives a two-dimensional computational space.
We choose the basis
\begin{equation}
  \begin{aligned}
    \ket{0_{\mathrm E}}&=
       \frac{\ket{\uparrow\downarrow\uparrow}
             -\ket{\downarrow\uparrow\uparrow}}{\sqrt{2}},\\
    \ket{1_{\mathrm E}}&=
       \frac{2\ket{\uparrow\uparrow\downarrow}
             -\ket{\uparrow\downarrow\uparrow}
             -\ket{\downarrow\uparrow\uparrow}}{\sqrt{6}}.
  \end{aligned}
  \label{eq:exchange-only-basis}
\end{equation}
The arrows refer to spins 1, 2, and 3.
Spins 1 and 2 form a singlet in $\ket{0_{\mathrm E}}$ and a triplet
in $\ket{1_{\mathrm E}}$, distinguishing the two $S=1/2$ representations.
The spin-rotation-invariant interaction derived in
Section~\ref{sec:exchange-virtual-tunneling} is called isotropic exchange.
It preserves $S,m$ but can mix these two states, allowing logical
rotations using exchange alone (Section~\ref{sec:encoded-gates}).

